\documentclass[10pt,twocolumn]{article}
\usepackage[scaled]{helvet}
\renewcommand\familydefault{\sfdefault}
\usepackage[T1]{fontenc}
\usepackage[margin=0.7in]{geometry}
\usepackage{titlesec}
\usepackage{enumitem}
\usepackage{parskip}
\setlength{\columnsep}{0.24in}
\setlist[itemize]{leftmargin=1.2em, topsep=0.2em, itemsep=0.18em}
\titleformat{\section}{\large\bfseries}{\thesection}{1em}{}

\begin{document}
\title{\vspace{-1.6cm}AWS Bedrock}
\author{AWS ML Study Notes}
\date{}
\maketitle
\small

\section*{Overview}
Bedrock is AWS's managed generative AI platform for working with foundation models through API based access without managing the underlying model serving infrastructure yourself.

It is useful when you want to build generative AI features on AWS while focusing on prompt design, retrieval, governance, and application behavior rather than GPU fleet operations.

\section*{Core Concepts}
\begin{itemize}
\item Bedrock provides access to multiple model families through a unified style of managed integration, along with features for guardrails, knowledge based augmentation, and agent oriented workflows.
\item The platform fits well with AWS identity, network, and observability controls, which can matter for enterprises that need private governance around generative AI usage.
\item The key abstraction is managed model invocation. You work with prompts, context, configurations, and outputs rather than building the serving stack yourself.
\end{itemize}

\section*{How It Works}
\begin{itemize}
\item An application sends a request containing the prompt and any supporting context to the chosen model endpoint through Bedrock's API.
\item Optional retrieval, policy, or orchestration layers can enrich the request before model inference, and the output is then returned to the application for postprocessing or user presentation.
\item Teams can wrap this with evaluation, prompt versioning, caching, and safety controls to create a governed production path.
\end{itemize}

\section*{AI and ML Relevance}
\begin{itemize}
\item Bedrock is relevant when ML teams need chat, summarization, extraction, or retrieval augmented generation features but do not want to self host large models.
\item It also enables a practical split of responsibilities: platform teams manage security and cost guardrails, while product teams iterate on prompts, context, and UX behavior.
\end{itemize}

\section*{Operational Notes}
\begin{itemize}
\item Do not confuse easy API access with solved product quality. Prompt design, retrieval quality, hallucination mitigation, and evaluation discipline still matter.
\item Watch token usage and latency, because generative AI cost scales differently from traditional endpoint hosting.
\item Treat private data carefully. Context sent to a model is still a data handling event and must respect governance boundaries.
\end{itemize}

\section*{What To Remember}
\begin{itemize}
\item Bedrock abstracts model hosting, not application design risk.
\item Retrieval quality and evaluation quality often matter more than swapping from one model to another.
\item For production genAI, governance and observability are as important as raw model capability.
\end{itemize}

\end{document}
